\documentclass[10pt,a4paper]{amsart}
\usepackage[margin=19mm]{geometry}
\usepackage{amsmath,amssymb,mathtools}
\usepackage[scr=rsfs,cal=euler]{mathalfa}
\usepackage{xfrac}
\usepackage[unicode,hidelinks,bookmarks=false]{hyperref}
\newcommand{\ie}{i.e.\ }
\newcommand{\cf}{cf.\ }
\newcommand{\resp}{respectively\ }
\newcommand{\Subsection}{\subsection}
\providecommand{\nobreakdash}{\nobreak\hbox{-}}
\setlength{\emergencystretch}{2em}
\multlinegap0pt
\usepackage{bm}
\usepackage{bbm}
\usepackage{dsfont}
\usepackage{xspace}
\usepackage{cancel}
\usepackage{mathtools}
\usepackage{amsthm}
\makeatletter
\@ifundefined{Theorem}{\newtheorem{Theorem}{Theorem}}{}
\@ifundefined{Lemma}{\newtheorem{Lemma}[Theorem]{Lemma}}{}
\@ifundefined{Remark}{\newtheorem{Remark}[Theorem]{Remark}}{}
\makeatother
\title[A Note on Pliability and the Openness of the Multiexponentia]{A Note on Pliability and the Openness of the Multiexponential Map in Carnot Groups}
\author{Fr\'ed\'eric Jean, Mario Sigalotti, Alessandro Socionovo}
\date{Source: arXiv:2507.13049v4 (CC BY 4.0)}
\begin{document}
\maketitle

\noindent\emph{Source and licence.} A Note on Pliability and the Openness of the Multiexponential Map in Carnot Groups. Version arXiv:2507.13049v4 (CC BY 4.0), distributed under the Creative Commons Attribution 4.0 International licence, as recorded in the arXiv licence metadata for this exact version and shown on the abstract page https://arxiv.org/abs/2507.13049v4. Full rights notice: licenses/arxiv-2507.13049-CC-BY-4.0.txt.\par\smallskip

\noindent\textit{Editorial scope.} Counted unit: the Lemma whose source label is \texttt{lem:strongly pliable implies weak-strong H} in the article (source0.tex lines 430--432, source label \texttt{lem:strongly pliable implies weak-strong H}), delivered together with the article's own complete original proof of it (source0.tex lines 434--470, one \texttt{proof} environment of 37 lines). No mathematics is written, repaired or re-derived: statement and proof are the article's own text line for line. Required context retained uncounted: eq:dilations-flow (lines 191--193); lem:pl\_iff\_strpl (lines 298--300); lem:pl\_iff\_strpl2 (lines 353--361); rmk:Vdense (lines 396--403). Every definition, display and label of those lines is retained unchanged. Omitted from this package and not redistributed: 1--190, 194--297, 301--352, 362--395, 404--429, 433--433, 471--597. Nothing inside a retained excerpt is omitted or altered: every retained body is delivered exactly as the article's lines are written, so no omission receipt of any kind accompanies this package. Citations: the retained excerpts carry no citation command at all, so no bibliography entry of the article is transcribed and no reference is redistributed. Reference adaptation: none was needed and none was made; every reference inside the delivered unit resolves inside it, so no unresolved reference or citation is produced. Printed slips kept literal: no printed word, symbol, hypothesis or number of the counted statement or of its proof was altered. Layout and class: the article's own class is \texttt{sigma} with packages including ifpdf, tikz, tikz-cd; the wrapper is the lane's pinned \texttt{amsart} editorial wrapper at 10pt on A4, which restates the theorem-like environments the retained text needs and adds the article's own macro definitions that the retained text needs (none), with extra wrapper packages (bm, bbm, dsfont, xspace, cancel, mathtools). The delivered numbering is the unit's own. Assets: no retained span contains a figure environment, a graphics-inclusion command, a PDF-page inclusion, a table or a TikZ picture; no image file is copied into this package, the package declares no asset dependency and no image file is redistributed with it. Licence: version arXiv:2507.13049v4 of this article is distributed under the Creative Commons Attribution 4.0 International licence, as recorded in the arXiv licence metadata for this exact version and shown on the abstract page https://arxiv.org/abs/2507.13049v4. A full rights notice is delivered at licenses/arxiv-2507.13049-CC-BY-4.0.txt. Characters: every retained excerpt is pure ASCII, so the delivered engine, XeLaTeX with its default Latin Modern text font, reports no missing character.
\begin{equation}\label{eq:dilations-flow}
\delta_\lambda (\exp(X)) = \exp (\delta_\lambda(X))
\end{equation}

\begin{Lemma}\label{lem:pl_iff_strpl}
 A vector $X\in{\mathfrak g}_1$ is pliable if and only if it is strongly pliable.
\end{Lemma}

\begin{Lemma} \label{lem:pl_iff_strpl2}
 Assume that
 $V$ is an
 $L^\infty$-stably $L^2$-dense subspace of
 $L^\infty(I,{\mathfrak g}_1)$.
 Then, a vector $X\in{\mathfrak g}_1$ is strongly
 pliable if and only if
 $X$ is strongly $V$-pliable.
\end{Lemma}

\begin{Remark}\label{rmk:Vdense}
 Let $V$ be
 an $L^\infty$-stably $L^2$-dense subspace of $ L^\infty(I,{\mathfrak g}_1)$.
 Then, by the two previous results, it follows that
 \begin{equation*}
 \text{str. $V$-pliable} \implies \text{$V$-pliable} \implies \text{pliable} \stackrel{\rm Lemma \ \ref{lem:pl_iff_strpl}}{\implies} \text{str. pliable} \stackrel{\rm Lemma \ \ref{lem:pl_iff_strpl2}}{\implies} \text{str. $V$-pliable}.
 \end{equation*}
\end{Remark}

\begin{Lemma}\label{lem:strongly pliable implies weak-strong H}
If $X\in {\mathfrak g}_1$ is strongly pliable, then it satisfies the {\rm (FH)}-condition.
\end{Lemma}

\begin{proof}
Let $X\in{\mathfrak g}_1$ be strongly pliable, and denote by ${\rm PC}={\rm PC}(I,{\mathfrak g}_1)$ the subspace of $L^\infty(I,{\mathfrak g}_1)$ made of piecewise constant functions with rational discontinuity points. By Remark~\ref{rmk:Vdense}, $X$ is strongly ${\rm PC}$-pliable.


 Fix $\eta>0$. By the strong ${\rm PC}$-pliability, there exist $Z,Y_1,\dots,Y_d\in {\rm PC}$, with $\|Z\|_\infty<\eta$, such that $E_X(Z) = E_X(0)$ and ${\rm d}E_X(Z)Y_1,\dots,{\rm d}E_X(Z)Y_d$ is a basis of $T_{E_X(0)}\mathbb{G}$. By definition of the space ${\rm PC}$, there exists an integer $p\ge 2$ such that \smash{$Z|_{((i-1)/p,i/p)}$}, \smash{$Y_1|_{((i-1)/p,i/p)},\dots,Y_{d}|_{((i-1)/p,i/p)}$} are constant for every $i=1,\dots,p$.

 Denote by $Z_i$ and $Y_{j,i}$ the values of $Z$ and $Y_j$ on the interval \smash{$\bigl(\frac{i-1}{p}, \frac ip\big)$}. Then the map $\Gamma^{(p)}\colon ({\mathfrak g}_1)^p\to\mathbb{G}$ defined by
 \[
 \Gamma^{(p)}(W_1,\dots,W_p):= {\rm e}^{{W_{p}}}\circ \dots \circ {\rm e}^{{W_{1}}}(1_\mathbb{G}),
 \]
satisfies $\Gamma^{(p)} (X+Z_1,\dots,X+Z_p) = \Gamma^{(p)} (X, \dots, X)$.
 It remains to show
that $\Gamma^{(p)}$ is a submersion at $(X+Z_1,\dots,X+Z_p)$.

We have that \smash{$Y_j=\sum_{i=1}^p Y_{j,i}\chi _{((i-1)/p,i/p)}$} and so
 \[{\rm d}E_X(Z)Y_j=\sum_{i=1}^p {\rm d}E_X(Z)Y_{j,i}\chi _{(\frac{i-1}p,\frac ip)}.\]
Moreover,
\begin{gather}
 {\rm d}E_X(Z)Y_{j,i}\chi _{(\frac{i-1}p,\frac ip)}\nonumber\\
 \qquad{}=\frac{{\rm d}}{{\rm d}\varepsilon}\biggr|_{\varepsilon=0}
 \Bigl({\rm e}^{\frac{X+Z_p}p}\circ\dots\circ {\rm e}^{\frac{X+Z_{i+1}}p}\circ {\rm e}^{\frac{X+Z_i+\varepsilon Y_{j,i}}p}\circ {\rm e}^{\frac{X+Z_{i-1}}p}\circ \dots\circ {\rm e}^{\frac{X+Z_1}p}(1_{\mathbb{G}})\Bigr)\nonumber\\
 \qquad{}\stackrel{\eqref{eq:dilations-flow}}{=}\frac{{\rm d}}{{\rm d}\varepsilon}\biggr|_{\varepsilon=0}
 \delta_{\frac1p}\bigl({\rm e}^{X+Z_p}\circ\dots\circ {\rm e}^{X+Z_{i+1}}\circ {\rm e}^{X+Z_i+\varepsilon Y_{j,i}}\circ {\rm e}^{X+Z_{i-1}}\circ \dots\circ {\rm e}^{X+Z_1}(1_{\mathbb{G}})\bigr)\nonumber\\
\qquad{}=\bigl(\delta_{\frac1p}\bigr)_*\frac{{\rm d}}{{\rm d}\varepsilon}\biggr|_{\varepsilon=0}\bigl({\rm e}^{X+Z_p}\circ\dots\circ {\rm e}^{X+Z_{i+1}}\circ {\rm e}^{X+Z_i+\varepsilon Y_{j,i}}\circ {\rm e}^{X+Z_{i-1}}\circ \dots\circ {\rm e}^{X+Z_1}(1_{\mathbb{G}})\bigr)\nonumber\\
 \qquad{}=\bigl(\delta_{\frac1p}\bigr)_*\frac{{\rm d}}{{\rm d}\varepsilon}\biggr|_{\varepsilon=0}
 \Gamma^{(p)}(X+Z_{1},\dots,X+Z_{i-1},X+Z_i+\varepsilon Y_{j,i},X+Z_{i+1}, \dots, X+Z_{p})\nonumber\\
\qquad{}=\bigl(\delta_{\frac1p}\bigr)_* {\rm d}\Gamma^{(p)}(X+Z_1,\dots,X+Z_p)\cdot (0,\dots,0, Y_{j,i},0,\dots,0).\label{eq:diffdilated}
\end{gather}
So
the image of the differential of
$\Gamma^{(p)}$
at $(X\!+\!Z_1,\dots,X\!+\!Z_p)$ contains the image by \smash{$(\delta_{p})_*\!=\!\bigl(\delta_{\frac1p}\bigr)_*^{-1}$} of
the
span of ${\rm d}E_X(Z)Y_1,\dots,{\rm d}E_X(Z)Y_d$.
Hence $\Gamma^{(p)}$
is a submersion at $(X\!+\!Z_1,\dots,X\!+\!Z_p)$.
\end{proof}

\end{document}
